% --- Archon physics formalization source begin ---
% archon:physics
% archon:covers PhyXMiniProblems/problem_phyx_mini_0720.lean
% archon:source-report reports/phyx_mini/problem_phyx_mini_0720.source.json

\chapter{Physics problem phyx\_mini\_0720}
\label{ch:PhyXMiniProblems_problem_phyx_mini_0720}

\paragraph{Problem source.}
\#\# Physical scenario

A \textbackslash{}(10.0\textbackslash{},\textbackslash{}mathrm\{cm\}\textbackslash{})-diameter suction cup is pushed against a smooth ceiling. The object can be suspended from the suction cup without pulling it off the ceiling? The mass of the suction cup is negligible.

\#\# Question

What is the maximum mass of an object?

\#\# Answer choices

- A: A: 77kg

- B: B: 79kg

- C: C: 81kg

- D: D: 83kg

\#\# Figure caption (auxiliary; use the image as primary evidence)

The image consists of two parts: a diagram on the left and a force diagram on the right.

Left Diagram:
- It shows an object hanging from a ceiling by a rod or string.
- The ceiling and object are represented in simplified forms, with the object labeled as "Object."

Right Force Diagram:
- It depicts a force diagram in the xy-plane with vectors:
  - \textbackslash{}( \textbackslash{}vec\{F\}\_\{\textbackslash{}text\{air\}\} \textbackslash{}) pointing upwards along the y-axis.
  - \textbackslash{}( \textbackslash{}vec\{F\}\_\{G\} \textbackslash{}) (Gravitational force) pointing downwards along the y-axis.
  - \textbackslash{}( \textbackslash{}vec\{F\}\_\{\textbackslash{}text\{net\}\} = \textbackslash{}vec\{0\} \textbackslash{}), indicating net force is zero, suggesting equilibrium.
  - A force labeled \textbackslash{}( \textbackslash{}vec\{n\} \textbackslash{}) (Normal force of ceiling) also pointing downwards.

Text labels indicate:
- "Normal force of ceiling" next to \textbackslash{}( \textbackslash{}vec\{n\} \textbackslash{}).
- "Gravitational force" next to \textbackslash{}( \textbackslash{}vec\{F\}\_\{G\} \textbackslash{}).

The combination of these elements visualizes forces acting on the object in equilibrium.

\#\# Dataset metadata

Statics; Physical Model Grounding Reasoning, Multi-Formula Reasoning

\paragraph{Recorded answer/context.}
C: C: 81kg

\paragraph{Figure/image path.}
/root/proposal\_for\_physic/science-mango/phyx\_mini\_run/phyx\_data/test\_image/720.png

\paragraph{Formalization target.}
create a compiling Lean file with sorry bodies at `PhyXMiniProblems/problem\_phyx\_mini\_0720.lean`. The Lean declarations must preserve the physical quantities, dimensions or dimensional roles, figure labels, governing-law hypotheses, and final relation expressed by this problem.
Use Mathlib/Physlib names found through LeanExplore where available. If a physics API is missing, introduce faithful local abstractions rather than scalar placeholder aliases.

\begin{theorem}[Physics formalization target]
\label{thm:physics:phyx_mini_0720:target}
\lean{PhyXMiniProblems.ProblemPhyXMini0720.problem_phyx_mini_0720}
\uses{def:physics:phyx-mini-0720:phyxminiproblems-problemphyxmini0720-massinkilograms, def:physics:phyx-mini-0720:phyxminiproblems-problemphyxmini0720-suctioncupsetup, def:physics:phyx-mini-0720:phyxminiproblems-problemphyxmini0720-matchessuctioncupscenario, def:physics:phyx-mini-0720:phyxminiproblems-problemphyxmini0720-matchessuppliedfigure, def:physics:phyx-mini-0720:phyxminiproblems-problemphyxmini0720-matchesproblemreadouts, def:physics:phyx-mini-0720:phyxminiproblems-problemphyxmini0720-usesstandardenvironment, def:physics:phyx-mini-0720:phyxminiproblems-problemphyxmini0720-usesidealvacuumlimit, def:physics:phyx-mini-0720:phyxminiproblems-problemphyxmini0720-satisfiescircularcupgeometry, def:physics:phyx-mini-0720:phyxminiproblems-problemphyxmini0720-satisfiessuctioncupstatics, def:physics:phyx-mini-0720:phyxminiproblems-problemphyxmini0720-atdetachmentthreshold, def:physics:phyx-mini-0720:phyxminiproblems-problemphyxmini0720-ismaximumsustainableobjectmass, def:physics:phyx-mini-0720:phyxminiproblems-problemphyxmini0720-isuniqueclosestanswer}
The assigned autoformalize agent should translate this physics problem into Lean declarations in the covered file, with theorem and lemma proof bodies written as `by sorry`.
\end{theorem}
\begin{proof}
This is an autoformalization task, not a proof task. Produce statements that can later be proved by the physics prover without weakening the physical model.
\end{proof}

% --- Archon declaration topology begin ---
\section{Declaration topology}
\begin{definition}[acceleration Dimension]
  \label{def:physics:phyx-mini-0720:phyxminiproblems-problemphyxmini0720-accelerationdimension}
  \lean{PhyXMiniProblems.ProblemPhyXMini0720.accelerationDimension}
  The physical dimension of an acceleration magnitude, L T⁻²
\end{definition}
\begin{proof}
  This is the declared model definition of acceleration Dimension.
\end{proof}
\begin{definition}[force Dimension]
  \label{def:physics:phyx-mini-0720:phyxminiproblems-problemphyxmini0720-forcedimension}
  \lean{PhyXMiniProblems.ProblemPhyXMini0720.forceDimension}
  \uses{def:physics:phyx-mini-0720:phyxminiproblems-problemphyxmini0720-accelerationdimension}
  The physical dimension of a force magnitude, M L T⁻²
\end{definition}
\begin{proof}
  This is the declared model definition of force Dimension.
\end{proof}
\begin{definition}[Length Quantity]
  \label{def:physics:phyx-mini-0720:phyxminiproblems-problemphyxmini0720-lengthquantity}
  \lean{PhyXMiniProblems.ProblemPhyXMini0720.LengthQuantity}
  A nonnegative unit-independent physical length
\end{definition}
\begin{proof}
  This is the declared model definition of Length Quantity.
\end{proof}
\begin{definition}[Mass Quantity]
  \label{def:physics:phyx-mini-0720:phyxminiproblems-problemphyxmini0720-massquantity}
  \lean{PhyXMiniProblems.ProblemPhyXMini0720.MassQuantity}
  A nonnegative unit-independent physical mass
\end{definition}
\begin{proof}
  This is the declared model definition of Mass Quantity.
\end{proof}
\begin{definition}[Acceleration Quantity]
  \label{def:physics:phyx-mini-0720:phyxminiproblems-problemphyxmini0720-accelerationquantity}
  \lean{PhyXMiniProblems.ProblemPhyXMini0720.AccelerationQuantity}
  \uses{def:physics:phyx-mini-0720:phyxminiproblems-problemphyxmini0720-accelerationdimension}
  A nonnegative unit-independent acceleration magnitude
\end{definition}
\begin{proof}
  This is the declared model definition of Acceleration Quantity.
\end{proof}
\begin{definition}[Force Magnitude]
  \label{def:physics:phyx-mini-0720:phyxminiproblems-problemphyxmini0720-forcemagnitude}
  \lean{PhyXMiniProblems.ProblemPhyXMini0720.ForceMagnitude}
  \uses{def:physics:phyx-mini-0720:phyxminiproblems-problemphyxmini0720-forcedimension}
  A nonnegative unit-independent force magnitude
\end{definition}
\begin{proof}
  This is the declared model definition of Force Magnitude.
\end{proof}
\begin{definition}[length In Meters]
  \label{def:physics:phyx-mini-0720:phyxminiproblems-problemphyxmini0720-lengthinmeters}
  \lean{PhyXMiniProblems.ProblemPhyXMini0720.lengthInMeters}
  \uses{def:physics:phyx-mini-0720:phyxminiproblems-problemphyxmini0720-lengthquantity}
  Coherent-SI metre readout of a physical length
\end{definition}
\begin{proof}
  This is the declared model definition of length In Meters.
\end{proof}
\begin{definition}[length In Centimeters]
  \label{def:physics:phyx-mini-0720:phyxminiproblems-problemphyxmini0720-lengthincentimeters}
  \lean{PhyXMiniProblems.ProblemPhyXMini0720.lengthInCentimeters}
  \uses{def:physics:phyx-mini-0720:phyxminiproblems-problemphyxmini0720-lengthquantity, def:physics:phyx-mini-0720:phyxminiproblems-problemphyxmini0720-lengthinmeters}
  Centimetre readout used for the stated cup diameter
\end{definition}
\begin{proof}
  This is the declared model definition of length In Centimeters.
\end{proof}
\begin{definition}[area In Square Meters]
  \label{def:physics:phyx-mini-0720:phyxminiproblems-problemphyxmini0720-areainsquaremeters}
  \lean{PhyXMiniProblems.ProblemPhyXMini0720.areaInSquareMeters}
  Coherent-SI square-metre readout of a physical area
\end{definition}
\begin{proof}
  This is the declared model definition of area In Square Meters.
\end{proof}
\begin{definition}[pressure In Pascals]
  \label{def:physics:phyx-mini-0720:phyxminiproblems-problemphyxmini0720-pressureinpascals}
  \lean{PhyXMiniProblems.ProblemPhyXMini0720.pressureInPascals}
  Coherent-SI pascal readout of a physical pressure
\end{definition}
\begin{proof}
  This is the declared model definition of pressure In Pascals.
\end{proof}
\begin{definition}[mass In Kilograms]
  \label{def:physics:phyx-mini-0720:phyxminiproblems-problemphyxmini0720-massinkilograms}
  \lean{PhyXMiniProblems.ProblemPhyXMini0720.massInKilograms}
  \uses{def:physics:phyx-mini-0720:phyxminiproblems-problemphyxmini0720-massquantity}
  Coherent-SI kilogram readout of a physical mass
\end{definition}
\begin{proof}
  This is the declared model definition of mass In Kilograms.
\end{proof}
\begin{definition}[acceleration In Meters Per Second Squared]
  \label{def:physics:phyx-mini-0720:phyxminiproblems-problemphyxmini0720-accelerationinmeterspersecondsquared}
  \lean{PhyXMiniProblems.ProblemPhyXMini0720.accelerationInMetersPerSecondSquared}
  \uses{def:physics:phyx-mini-0720:phyxminiproblems-problemphyxmini0720-accelerationquantity}
  Metre-per-second-squared readout of an acceleration magnitude
\end{definition}
\begin{proof}
  This is the declared model definition of acceleration In Meters Per Second Squared.
\end{proof}
\begin{definition}[force In Newtons]
  \label{def:physics:phyx-mini-0720:phyxminiproblems-problemphyxmini0720-forceinnewtons}
  \lean{PhyXMiniProblems.ProblemPhyXMini0720.forceInNewtons}
  \uses{def:physics:phyx-mini-0720:phyxminiproblems-problemphyxmini0720-forcemagnitude}
  Coherent-SI newton readout of a force magnitude
\end{definition}
\begin{proof}
  This is the declared model definition of force In Newtons.
\end{proof}
\begin{definition}[pressure Difference Force In Newtons]
  \label{def:physics:phyx-mini-0720:phyxminiproblems-problemphyxmini0720-pressuredifferenceforceinnewtons}
  \lean{PhyXMiniProblems.ProblemPhyXMini0720.pressureDifferenceForceInNewtons}
  \uses{def:physics:phyx-mini-0720:phyxminiproblems-problemphyxmini0720-areainsquaremeters, def:physics:phyx-mini-0720:phyxminiproblems-problemphyxmini0720-pressureinpascals}
  The upward force in newtons supplied by a uniform pressure difference across the sealed circular area
\end{definition}
\begin{proof}
  This is the declared model definition of pressure Difference Force In Newtons.
\end{proof}
\begin{definition}[weight In Newtons]
  \label{def:physics:phyx-mini-0720:phyxminiproblems-problemphyxmini0720-weightinnewtons}
  \lean{PhyXMiniProblems.ProblemPhyXMini0720.weightInNewtons}
  \uses{def:physics:phyx-mini-0720:phyxminiproblems-problemphyxmini0720-massquantity, def:physics:phyx-mini-0720:phyxminiproblems-problemphyxmini0720-accelerationquantity, def:physics:phyx-mini-0720:phyxminiproblems-problemphyxmini0720-massinkilograms, def:physics:phyx-mini-0720:phyxminiproblems-problemphyxmini0720-accelerationinmeterspersecondsquared}
  The weight magnitude in newtons of a physical mass
\end{definition}
\begin{proof}
  This is the declared model definition of weight In Newtons.
\end{proof}
\begin{definition}[Vertical Direction]
  \label{def:physics:phyx-mini-0720:phyxminiproblems-problemphyxmini0720-verticaldirection}
  \lean{PhyXMiniProblems.ProblemPhyXMini0720.VerticalDirection}
  Vertical directions of the three nonzero arrows in the supplied diagram
\end{definition}
\begin{proof}
  This is the declared model definition of Vertical Direction.
\end{proof}
\begin{definition}[Figure Force Label]
  \label{def:physics:phyx-mini-0720:phyxminiproblems-problemphyxmini0720-figureforcelabel}
  \lean{PhyXMiniProblems.ProblemPhyXMini0720.FigureForceLabel}
  Force labels printed in the supplied raster
\end{definition}
\begin{proof}
  This is the declared model definition of Figure Force Label.
\end{proof}
\begin{definition}[Suction Cup Figure]
  \label{def:physics:phyx-mini-0720:phyxminiproblems-problemphyxmini0720-suctioncupfigure}
  \lean{PhyXMiniProblems.ProblemPhyXMini0720.SuctionCupFigure}
  \uses{def:physics:phyx-mini-0720:phyxminiproblems-problemphyxmini0720-verticaldirection, def:physics:phyx-mini-0720:phyxminiproblems-problemphyxmini0720-figureforcelabel}
  Qualitative and textual evidence from the two-part supplied figure. This record contains no pressure, mass, or force magnitude
\end{definition}
\begin{proof}
  This is the declared model definition of Suction Cup Figure.
\end{proof}
\begin{definition}[Attachment State]
  \label{def:physics:phyx-mini-0720:phyxminiproblems-problemphyxmini0720-attachmentstate}
  \lean{PhyXMiniProblems.ProblemPhyXMini0720.AttachmentState}
  Mechanical state of the suction cup and its suspended load
\end{definition}
\begin{proof}
  This is the declared model definition of Attachment State.
\end{proof}
\begin{definition}[Suction Cup Setup]
  \label{def:physics:phyx-mini-0720:phyxminiproblems-problemphyxmini0720-suctioncupsetup}
  \lean{PhyXMiniProblems.ProblemPhyXMini0720.SuctionCupSetup}
  \uses{def:physics:phyx-mini-0720:phyxminiproblems-problemphyxmini0720-lengthquantity, def:physics:phyx-mini-0720:phyxminiproblems-problemphyxmini0720-massquantity, def:physics:phyx-mini-0720:phyxminiproblems-problemphyxmini0720-accelerationquantity, def:physics:phyx-mini-0720:phyxminiproblems-problemphyxmini0720-forcemagnitude, def:physics:phyx-mini-0720:phyxminiproblems-problemphyxmini0720-suctioncupfigure, def:physics:phyx-mini-0720:phyxminiproblems-problemphyxmini0720-attachmentstate}
  Independent physical quantities of the setup. The suspended mass and the three force magnitudes are observables constrained only by the hypotheses below; they are not defined from the recorded answer
\end{definition}
\begin{proof}
  This is the declared model definition of Suction Cup Setup.
\end{proof}
\begin{definition}[Matches Suction Cup Scenario]
  \label{def:physics:phyx-mini-0720:phyxminiproblems-problemphyxmini0720-matchessuctioncupscenario}
  \lean{PhyXMiniProblems.ProblemPhyXMini0720.MatchesSuctionCupScenario}
  \uses{def:physics:phyx-mini-0720:phyxminiproblems-problemphyxmini0720-suctioncupsetup}
  The qualitative apparatus assumptions stated in the problem
\end{definition}
\begin{proof}
  This is the declared model definition of Matches Suction Cup Scenario.
\end{proof}
\begin{definition}[Matches Supplied Figure]
  \label{def:physics:phyx-mini-0720:phyxminiproblems-problemphyxmini0720-matchessuppliedfigure}
  \lean{PhyXMiniProblems.ProblemPhyXMini0720.MatchesSuppliedFigure}
  \uses{def:physics:phyx-mini-0720:phyxminiproblems-problemphyxmini0720-suctioncupsetup}
  Primary-image evidence: F\_air points upward, F\_G and n point downward, the net force is shown as zero, and the object hangs beneath the ceiling cup
\end{definition}
\begin{proof}
  This is the declared model definition of Matches Supplied Figure.
\end{proof}
\begin{definition}[Matches Problem Readouts]
  \label{def:physics:phyx-mini-0720:phyxminiproblems-problemphyxmini0720-matchesproblemreadouts}
  \lean{PhyXMiniProblems.ProblemPhyXMini0720.MatchesProblemReadouts}
  \uses{def:physics:phyx-mini-0720:phyxminiproblems-problemphyxmini0720-lengthincentimeters, def:physics:phyx-mini-0720:phyxminiproblems-problemphyxmini0720-massinkilograms, def:physics:phyx-mini-0720:phyxminiproblems-problemphyxmini0720-suctioncupsetup}
  Numerical and idealizing data stated or implied by the problem. The suction cup diameter is 10.0 cm, and its negligible mass is modeled as zero. This record contains no suspended-object mass or answer-choice selection
\end{definition}
\begin{proof}
  This is the declared model definition of Matches Problem Readouts.
\end{proof}
\begin{definition}[Uses Standard Environment]
  \label{def:physics:phyx-mini-0720:phyxminiproblems-problemphyxmini0720-usesstandardenvironment}
  \lean{PhyXMiniProblems.ProblemPhyXMini0720.UsesStandardEnvironment}
  \uses{def:physics:phyx-mini-0720:phyxminiproblems-problemphyxmini0720-pressureinpascals, def:physics:phyx-mini-0720:phyxminiproblems-problemphyxmini0720-accelerationinmeterspersecondsquared, def:physics:phyx-mini-0720:phyxminiproblems-problemphyxmini0720-suctioncupsetup}
  Standard environmental calibration used by the recorded numerical answer: one standard atmosphere is 101325 Pa and near-Earth gravity is 9.8 m/s²
\end{definition}
\begin{proof}
  This is the declared model definition of Uses Standard Environment.
\end{proof}
\begin{definition}[Uses Ideal Vacuum Limit]
  \label{def:physics:phyx-mini-0720:phyxminiproblems-problemphyxmini0720-usesidealvacuumlimit}
  \lean{PhyXMiniProblems.ProblemPhyXMini0720.UsesIdealVacuumLimit}
  \uses{def:physics:phyx-mini-0720:phyxminiproblems-problemphyxmini0720-pressureinpascals, def:physics:phyx-mini-0720:phyxminiproblems-problemphyxmini0720-suctioncupsetup}
  The maximum ideal suction is attained when the pressure beneath the sealed cup is vacuum. This is a pressure calibration, not a mass conclusion
\end{definition}
\begin{proof}
  This is the declared model definition of Uses Ideal Vacuum Limit.
\end{proof}
\begin{definition}[Satisfies Circular Cup Geometry]
  \label{def:physics:phyx-mini-0720:phyxminiproblems-problemphyxmini0720-satisfiescircularcupgeometry}
  \lean{PhyXMiniProblems.ProblemPhyXMini0720.SatisfiesCircularCupGeometry}
  \uses{def:physics:phyx-mini-0720:phyxminiproblems-problemphyxmini0720-lengthinmeters, def:physics:phyx-mini-0720:phyxminiproblems-problemphyxmini0720-areainsquaremeters, def:physics:phyx-mini-0720:phyxminiproblems-problemphyxmini0720-suctioncupsetup}
  The sealed face is the disk determined by the stated cup diameter
\end{definition}
\begin{proof}
  This is the declared model definition of Satisfies Circular Cup Geometry.
\end{proof}
\begin{definition}[Satisfies Suction Cup Statics]
  \label{def:physics:phyx-mini-0720:phyxminiproblems-problemphyxmini0720-satisfiessuctioncupstatics}
  \lean{PhyXMiniProblems.ProblemPhyXMini0720.SatisfiesSuctionCupStatics}
  \uses{def:physics:phyx-mini-0720:phyxminiproblems-problemphyxmini0720-forceinnewtons, def:physics:phyx-mini-0720:phyxminiproblems-problemphyxmini0720-pressuredifferenceforceinnewtons, def:physics:phyx-mini-0720:phyxminiproblems-problemphyxmini0720-weightinnewtons, def:physics:phyx-mini-0720:phyxminiproblems-problemphyxmini0720-suctioncupsetup}
  Governing pressure-force and weight laws, together with vertical static equilibrium. The diagram's upward air force balances the downward ceiling normal and gravitational forces. No numerical mass occurs in these laws
\end{definition}
\begin{proof}
  This is the declared model definition of Satisfies Suction Cup Statics.
\end{proof}
\begin{definition}[At Detachment Threshold]
  \label{def:physics:phyx-mini-0720:phyxminiproblems-problemphyxmini0720-atdetachmentthreshold}
  \lean{PhyXMiniProblems.ProblemPhyXMini0720.AtDetachmentThreshold}
  \uses{def:physics:phyx-mini-0720:phyxminiproblems-problemphyxmini0720-forceinnewtons, def:physics:phyx-mini-0720:phyxminiproblems-problemphyxmini0720-suctioncupsetup}
  At incipient detachment the ceiling can no longer exert a downward normal force on the cup, so that force has fallen to zero
\end{definition}
\begin{proof}
  This is the declared model definition of At Detachment Threshold.
\end{proof}
\begin{definition}[Can Be Suspended]
  \label{def:physics:phyx-mini-0720:phyxminiproblems-problemphyxmini0720-canbesuspended}
  \lean{PhyXMiniProblems.ProblemPhyXMini0720.CanBeSuspended}
  \uses{def:physics:phyx-mini-0720:phyxminiproblems-problemphyxmini0720-massquantity, def:physics:phyx-mini-0720:phyxminiproblems-problemphyxmini0720-pressuredifferenceforceinnewtons, def:physics:phyx-mini-0720:phyxminiproblems-problemphyxmini0720-weightinnewtons, def:physics:phyx-mini-0720:phyxminiproblems-problemphyxmini0720-suctioncupsetup}
  A candidate object mass is sustainable when its weight plus the negligible cup weight does not exceed the available pressure-difference force
\end{definition}
\begin{proof}
  This is the declared model definition of Can Be Suspended.
\end{proof}
\begin{definition}[Is Maximum Sustainable Object Mass]
  \label{def:physics:phyx-mini-0720:phyxminiproblems-problemphyxmini0720-ismaximumsustainableobjectmass}
  \lean{PhyXMiniProblems.ProblemPhyXMini0720.IsMaximumSustainableObjectMass}
  \uses{def:physics:phyx-mini-0720:phyxminiproblems-problemphyxmini0720-massquantity, def:physics:phyx-mini-0720:phyxminiproblems-problemphyxmini0720-massinkilograms, def:physics:phyx-mini-0720:phyxminiproblems-problemphyxmini0720-suctioncupsetup, def:physics:phyx-mini-0720:phyxminiproblems-problemphyxmini0720-canbesuspended}
  A mass is maximal among all object masses sustainable by this setup
\end{definition}
\begin{proof}
  This is the declared model definition of Is Maximum Sustainable Object Mass.
\end{proof}
\begin{definition}[Answer Choice]
  \label{def:physics:phyx-mini-0720:phyxminiproblems-problemphyxmini0720-answerchoice}
  \lean{PhyXMiniProblems.ProblemPhyXMini0720.AnswerChoice}
  Labels of the four mass choices printed with the problem
\end{definition}
\begin{proof}
  This is the declared model definition of Answer Choice.
\end{proof}
\begin{definition}[displayed Mass In Kilograms]
  \label{def:physics:phyx-mini-0720:phyxminiproblems-problemphyxmini0720-displayedmassinkilograms}
  \lean{PhyXMiniProblems.ProblemPhyXMini0720.displayedMassInKilograms}
  \uses{def:physics:phyx-mini-0720:phyxminiproblems-problemphyxmini0720-answerchoice}
  Kilogram values printed beside the answer labels
\end{definition}
\begin{proof}
  This is the declared model definition of displayed Mass In Kilograms.
\end{proof}
\begin{definition}[Is Unique Closest Answer]
  \label{def:physics:phyx-mini-0720:phyxminiproblems-problemphyxmini0720-isuniqueclosestanswer}
  \lean{PhyXMiniProblems.ProblemPhyXMini0720.IsUniqueClosestAnswer}
  \uses{def:physics:phyx-mini-0720:phyxminiproblems-problemphyxmini0720-massinkilograms, def:physics:phyx-mini-0720:phyxminiproblems-problemphyxmini0720-suctioncupsetup, def:physics:phyx-mini-0720:phyxminiproblems-problemphyxmini0720-answerchoice, def:physics:phyx-mini-0720:phyxminiproblems-problemphyxmini0720-displayedmassinkilograms}
  A displayed choice is uniquely closest to the derived limiting mass
\end{definition}
\begin{proof}
  This is the declared model definition of Is Unique Closest Answer.
\end{proof}
% --- Archon declaration topology end ---
% --- Archon physics formalization source end ---
